\clearpage
\section{Reading an actual lease transition}
\label{sec:pingan-lease-transition}

Ping An's 2019 interim report shows what happens when the dated lease rules
meet an actual disclosure. The group reports adoption of revised ASBE~21 on
1 January 2019 and leaves its 2018 comparative financial statements
unadjusted. Its stated listings in Shanghai and Hong Kong place it within
the enterprise class covered by the 2018 notice's 2019 application date
\citep[pp.\ 112--113]{pingan2019asbe};
\citep[notice, para.\ I]{moflease2018}. This supplies an actual adoption
fact for the representation. The later deferred-tax treatment described
above still requires its own applicable date and tax facts.

The immediate reconstruction problem is to keep three observations apart:
the original December presentation, the January adoption balances and the
June closing balances. A right-of-use asset records the group's recognised
right to use leased property; the lease liability measures the associated
recognised payment obligation. Table~\ref{tab:pingan-lease-balances} retains
the amounts and their dates. All amounts in this example are displayed
millions of renminbi; arithmetic on them does not establish exact underlying
amounts or the issuer's rounding rule.

\begin{table}[!ht]
\centering
\caption{Adoption balances and later carrying balances answer different questions.}
\label{tab:pingan-lease-balances}
\begin{tabular}{@{}lrr@{}}
\toprule
Observation & Right-of-use assets & Lease liabilities \\
\midrule
31 December 2018, unadjusted comparative & --- & --- \\
1 January 2019, reported at adoption & 18,500 & 17,667 \\
30 June 2019, carrying balance & 17,390 & 17,212 \\
\bottomrule
\end{tabular}
\end{table}

The comparative dashes reproduce the presentation under the previous basis;
they do not say that the group had no lease contracts in December. The
January amounts come from the policy-change note, whereas the June amounts
come from the consolidated balance sheet
\citep[pp.\ 99--100, 112]{pingan2019asbe}. Subtracting a December dash
treated as zero from the June liability would mix recognition at transition
with six months of subsequent activity.

The opening asset exceeds the liability by $18{,}500-17{,}667=833$. That
difference alone does not identify an increase in retained earnings.
Article~61(2) permits specified alternative measurements of the right-of-use
asset for former operating leases, including a liability-based amount
adjusted for prepaid rent; its cumulative approach adjusts retained earnings
and other relevant financial-statement items
\citep[attachment, art.\ 61(2)]{moflease2018}. To see why the other items
matter, suppose solely for illustration that pre-existing prepaid rent of
833 is reclassified. Debiting the new right-of-use asset by 18,500 and
crediting the new liability by 17,667 and prepaid rent by 833 balances the
entry with no equity movement. Assigning the same residual entirely to
retained earnings would instead require a different assumption about all
other transition adjustments. The report does not identify the illustrative
prepayment or supply a complete opening journal. Reconstruction therefore
retains the unexplained components rather than creating an equity plug.

The policy note nevertheless narrows the permitted interpretation. It says
the opening liability discounts remaining payments at the lessee's
incremental borrowing rate on 1 January 2019. It names five simplifications:
treating operating leases with a remaining term shorter than twelve months
as short-term leases; excluding initial direct costs at transition; using
one discount rate for leases with similar characteristics; using actual
option exercise and other recent information to determine the lease term;
and assessing pre-existing onerous contracts instead of a separate
right-of-use impairment test. These correspond to the specified provisions
of Article~63, with the rate and direct-cost choices appearing in the same
numbered provision \citep[p.\ 113]{pingan2019asbe};
\citep[attachment, art.\ 63]{moflease2018}. The record should preserve
the reported elections without assigning an unreported rate or deciding
which asset-measurement alternative was used for each lease. The note also
calls the difference between discounted prior operating-lease commitments
and the recognised opening liability ``not material''. That qualification
supplies neither an equality nor a numeric tolerance.

The June maturity disclosure adds a different measurement of the obligation.
Under its heading for undiscounted contractual and expected cash flows,
the lease-liability row gives 1,515 within three months, 5,885 in three to
twelve months, 11,048 in one to five years and 485 beyond five years. Those
numeric cells sum to 18,933; the immediate and no-term columns contain
dashes \citep[pp.\ 164--165]{pingan2019asbe}. The first two dated buckets
sum to 7,400, but disclose no payment dates within those buckets. The
18,933 total and the 17,212 carrying liability can coexist because they
measure different things. Their difference, 1,721, is not the interest
expense for the half-year and does not by itself determine a discount rate.

Even on the post-adoption basis, a change in a balance is not a complete
account of the transactions that produced it. The right-of-use asset falls
by $17{,}390-18{,}500=-1{,}110$ and the liability by
$17{,}212-17{,}667=-455$. Additions, remeasurements, disposals and other
movements can accompany depreciation or payments. The first difference
therefore does not identify depreciation, and the second does not identify
cash paid. The consolidated cash-flow statement's finance-leasing business
borrowing and receivable lines concern the group's financing activities as
a business; their labels do not make them payments on its own lessee
obligations \citep[pp.\ 106--107]{pingan2019asbe}.

The same care prevents a false narrative conflict. In its major-contract
discussion, the company says that it has no reportable custody, contracting,
lease or other major-contract matters for the period
\citep[p.\ 86]{pingan2019asbe}. Preserving the company reference and
the reportable-major-contract qualification gives a different assertion
from saying that the consolidated group has no leases. Removing those
qualifiers would create the apparent contradiction with the balance sheet.

For the proposed representation, the useful result is a set of dated,
source-supported balances, policy elections and maturity restrictions.
A query for the reported June carrying liability can return 17,212 with
its basis and precision. A query for future scheduled payments must keep
the disclosed maturity buckets and leave their internal dates unresolved.
A query for the complete transition journal or the half-year's lessee cash
payments must retain the missing inputs. These distinct answers let actual
disclosures constrain histories without pretending that every balanced
completion is the company's history.

\clearpage
